\documentclass[12pt,a4paper]{article}
\usepackage[brazilian]{babel}
\usepackage[utf8]{inputenc}
\usepackage{amsmath,amssymb,amsthm}
\usepackage{geometry}
\usepackage{enumitem}
\usepackage{hyperref}
\usepackage{xcolor}
\usepackage{graphicx}
\usepackage{fancyhdr}
\usepackage{tcolorbox}
\usepackage{totpages}
\usepackage{titlesec}
\usepackage{etoolbox}
\usepackage{multicol}
\usepackage{multicol}
\setlength{\columnseprule}{0.4pt} % Espessura da linha divisória
\renewcommand{\columnseprulecolor}{\color{borderlight}} % Cor da linha (opcional)

\geometry{top=3.5cm, bottom=2.5cm, left=2.0cm, right=2.0cm, headheight=2.6cm}

\definecolor{primary}{RGB}{10, 45, 90}
\definecolor{secondary}{RGB}{25, 110, 180}
\definecolor{bglight}{RGB}{245, 247, 250}
\definecolor{borderlight}{RGB}{210, 220, 230}
\definecolor{correctgreen}{RGB}{40, 140, 60}

\tcbuselibrary{skins, breakable}

\newtcolorbox{boxfixacao}[1]{
  enhanced,
  colback=bglight,
  colframe=primary,
  arc=3pt,
  boxrule=1pt,
  title={#1},
  fonttitle=\sffamily\bfseries\color{white},
  coltitle=white,
  colbacktitle=primary,
  breakable
}

\newcounter{numquestao}
\setcounter{numquestao}{0}

\newenvironment{questao}[2][]{%
  \refstepcounter{numquestao}%
  \vspace{0.3cm}%
  \begin{tcolorbox}[
    enhanced,
    colback=white,
    colframe=secondary,
    arc=2pt,
    boxrule=0.8pt,
    leftrule=4pt,
    title={\textbf{Questão \thenumquestao} \ifstrempty{#2}{}{--- #2} \hfill \small\textnormal{#1}},
    fonttitle=\sffamily\bfseries\color{white},
    coltitle=white,
    colbacktitle=secondary,
    breakable
  ]%
}{%
  \end{tcolorbox}%
}

\newenvironment{gabarito}{%
  \vspace{0.1cm}%
  \begin{tcolorbox}[
    enhanced,
    colback=bglight,
    colframe=correctgreen,
    arc=2pt,
    boxrule=0.6pt,
    title={\small\sffamily\bfseries Resolução / Gabarito},
    fonttitle=\sffamily\color{white},
    colbacktitle=correctgreen,
    breakable
  ]%
  \small
}{%
  \end{tcolorbox}%
}

\newcommand{\espacoresposta}[1]{%
  \vspace{#1}%
}

\pagestyle{fancy}
\fancyhf{}

\fancyhead[L]{%
  \small\sffamily
  \textbf{IFPE -- Campus Caruaru}
}
\fancyhead[R]{%
  \small\sffamily
  \textbf{Instituto Federal de Pernambuco -- Campus Caruaru} \\
  \textcolor{secondary}{Disciplina:} Cálculo 1 \\
  \textcolor{secondary}{Prof.:} Cleibson José da Silva
}

\fancyfoot[L]{\small\sffamily\color{gray}Lista 6 -- Integrais por Substituição}
\fancyfoot[R]{\small\sffamily\color{gray}Página \thepage\ de \ref{TotPages}}

\renewcommand{\headrulewidth}{0.8pt}
\renewcommand{\headrule}{{\color{primary}\hrule width\headwidth height \headrulewidth}}
\renewcommand{\footrulewidth}{0.4pt}

\hypersetup{
    colorlinks=true,
    linkcolor=secondary,
    urlcolor=secondary,
    citecolor=primary
}

\begin{document}


\vspace{0.2cm}

\begin{center}
  {\LARGE\bfseries\color{primary} Lista 6: Integrais por substituição} \\[0.1cm]
  {\small\sffamily\color{gray} Exercícios sobre integrais por substituição}
\end{center}




\vspace{0.4cm}

\begin{enumerate}
 \item Nas questões abaixo calcule a integral fazendo a substituição dada.
 
 
 \begin{enumerate}

\item $\displaystyle \int x(4 + x^2)^{10} \, dx, \quad u = 4 + x^2$

\item $\displaystyle \int \frac{dt}{(1 - 6t)^4}, \quad u = 1 - 6t$

\item $\displaystyle \int e^{\operatorname{sen} \theta} \cos \theta \, d\theta, \quad u = \operatorname{sen} \theta$

\end{enumerate}

\item Calcule a integral indefinida.

\begin{enumerate}

    \item $\displaystyle \int x^2(x^3 + 5)^9 \, dx$

    \item $\displaystyle \int (3t + 2)^{2.4} \, dt$

    \item $\displaystyle \int \frac{x}{(x^2 + 1)^2} \, dx$

    \item $\displaystyle \int e^x \operatorname{sen}(e^x) \, dx$

    \item $\displaystyle \int \frac{x}{x^2 + 1} \, dx$

    \item $\displaystyle \int \sec 2\theta \operatorname{tg} 2\theta \, d\theta$

    \item $\displaystyle \int \frac{dx}{ax + b} \quad (a \neq 0)$

    \item $\displaystyle \int \sqrt{x} \operatorname{sen}(1 + x^{3/2}) \, dx$

    \item $\displaystyle \int (1 + \operatorname{tg} \theta)^5 \sec^2 \theta \, d\theta$

    \item $\displaystyle \int e^{\cos t} \operatorname{sen} t \, dt$
\end{enumerate}

\end{enumerate}


\end{document}